\section{Evaluation}
\label{sec:eval}
%% Structure (keep): RQ list -> Subjects -> Baselines -> Measurements
%% -> Setup -> one \subsection per RQ with Observations.
%% Each RQ is a plain question a practitioner would ask.

\focus{\textbf{Research questions.} State 2--4 questions a practitioner would
  actually ask. Each is a plain yes/no or how-much question you can answer
  with a measurement --- not ``Is our tool good?''. Cover: does it work
  (\rqref{1}), is it better than a simpler baseline (\rqref{2}), does it
  change the outcome for a user (\rqref{3}), and what does it cost
  (\rqref{4}). Every RQ must be answered later by an \texttt{observation}
  box backed by a number. This is graded.}

We ask \todo{n} questions.
\begin{itemize}
  \item[\rqref{1}] \todo{Question?}
  \item[\rqref{2}] \todo{Question, compared with simpler baselines?}
  \item[\rqref{3}] \todo{Does it change what tools/users get?}
  \item[\rqref{4}] What does \tool cost in time \todo{and resources}?
\end{itemize}

\subsection{Subjects}
\label{sec:subjects}
%% Selection criteria as an enumerate; exact versions/commits; why these.
We select \todo{n} subjects with the following criteria:
\begin{enumerate}
  \item \todo{criterion, with the reason};
\end{enumerate}

\subsection{Baselines}
\label{sec:baselines}
%% description list; each baseline: what it is, why it is a fair comparison.
\begin{description}
  \item[Default.] \todo{}
\end{description}

\subsection{Measurements}
\begin{description}
  \item[\rqref{1}.] \todo{What exactly is measured.}
\end{description}

\subsection{Setup}
%% Machine, OS, compiler/tool versions, time limits, repetitions.
\todo{Machine, versions, repetitions.}

\subsection{\rqref{1}: \todo{Short question}}
\label{sec:rq1}
\todo{Table/figure first, then text that reads it.}

\begin{observation}
\label{obs:first}
\todo{One claim, with the number that supports it.}
\end{observation}
